\documentclass{syllabus}

\begin{document}

\thispagestyle{fancy}
\date{Fall 2025}
\author{}
\title{PHY 1200: Fundamental Physics \\ {\large \mdseries An introduction to physical principles and a survey of disciplines}}
\maketitle



\section*{General Information}

\raisebox{0.9\height}{
 % TODO: Change this
\begin{tabular}{lll}
\textbf{Instructor:}    & \multicolumn{2}{l}{Joseph Anderson}                             \\
\textbf{}               & \multicolumn{2}{l}{\href{mailto:janderson19@carthage.edu}{janderson19@carthage.edu}}                    \\
\textbf{}               & \multicolumn{2}{l}{DSC 082}                    \\
\textbf{}               & \multicolumn{2}{l}{(262)551-5871}                               \\ \\
\textbf{Office hours:}  & Mon, Wed               & 10:20 AM - 12:00 PM                    \\ \\
\textbf{Textbook:}      & \multicolumn{2}{l}{\textit{University Physics}} \\
\textbf{}               & \multicolumn{2}{l}{Open Stax}     \\ \\
\textbf{Prerequisites:} & \multicolumn{2}{l}{Co-enrolled in one of the following}                    \\
\textbf{}               & \multicolumn{2}{l}{MTH 1070 (functions and graphs)}  \\
\textbf{}               & \multicolumn{2}{l}{MTH 1120 or 1220 (Calculus I or II)} 
\end{tabular}
%
}
\hspace{\fill}
\includegraphics[height=2.25in]{Assets/openstax.png}
\hspace{.25in}

\section*{Course overview}
This course covers fundamental physical principles including descriptions of mechanical, electrical, wave, and atomic phenomena. The course highlights ways in which physical principles are used to describe and understand the vast array of observable phenomena in the universe. Students will study applications of physics to a range of important historical and contemporary scientific and technological questions. This course is intended for potential physics majors or students planning further study in the physical sciences.
% TODO: add my own description

The purpose of this course is to whet your appetite for the rest of the physics curriculum; it provides a survey not only of the various topics that physics studies, but the types of principles used in each of these disciplines of physics. You will also learn how to look at math from an applied perspective and the habits of thought and work needed to succeed in physics. 


\subsection*{Student learning outcomes} 

This course is structured as an environment to cultivate the following personal traits:
\begin{enumerate}
    \item \textbf{Understanding:} the personal quality by which one sees particular facts in light of underlying principles.
    \item \textbf{Curiosity:} the disposition toward asking questions about the world and seeks to understand why things are the way they are.
\end{enumerate}
Although the course will be organized so as to encourage this development, you are the only one who can determine whether to do the developing or not. Nonetheless, course assessment will be primarily structured around the following outcomes. 


Upon successful completion of this course, the student will be able to:
\begin{enumerate}[label=\textcolor{CarthageDark}{\bfseries\arabic*.}]
    \item Recognize quantifiable relationships between physical quantities and suggest appropriate mathematical forms for these relationships.
    \item Compare and contrast the various fields of physics as well as the types of principles/methods used in each.
    \item Outline the major mathematical principles commonly used in physics, including vector algebra, differential reasoning, and other geometric concepts.
    \item Explain the role of both physical experimentation and `thought experiments' in arriving at physical principles.
    \item Apply mathematics and physical principles to make calculations and predict the outcome of physical scenarios.
\end{enumerate}

\section*{Course components}
\subsection{Physics field work}
Science involves being attentive to causal relationships in the physical world. These assignments will involve noticing a phenomenon in daily experience and analyzing it from the perspective of physics. On a weekly basis, you will be asked to identify an everyday phenomenon that exemplifies either: \textbf{1)} a physical principle from class or \textbf{2)} a mathematical principle from class. 

\subsection{Homework}
Homework in this class is designed for you to get used to the types of problem-solving done in physics. Homework assignments in this course are purposefully short and will have a mixture of mathematical calculations and problem-solving using physical relationships to determine some unknown quantity. 

\subsection{Discussion reflection}
Part of this course will involve reflection on questions relating to the foundations and practice of science. These topics will usually be accompanied by class discussions, with a half-page reflection on the topic to follow which will be used as a measure of your participation in the class discussion. 

\subsection{Campus engagement}
Lone wolves do not succeed in anything, and physics is no exception. Engaging in scientific communities is a skill as well as an important aid in your own academic success. A small portion of your grade will be based on engagement with the society of physics students (SPS) (or another physics related club) and other campus events. 

For full credit on this portion of your grade, you must attend seven campus events, four (4) of which must be meetings of the SPS. These may be regular meetings (monthly), social events (sporadic), or homework hang-out events (weekly). If you are not a physics major, you may suggest clubs in your own discipline (e.g. Chem club) to fulfill this requirement, but these substitutions must be approved by the instructor before the first event credit is due. 

\subsection{Laboratory activities}
These will involve a worksheet to be completed during in the course of the lab. These will usually be able to be completed in the laboratory period, but can be finished the night of, to be due during the first class session following the lab period. 

A small portion of your lab grade will be an exit ticket reflection on the role of physical principles in the laboratory activity.

\subsection{Exams}
Exams are designed to stretch your understanding to test how you can apply the principles learned in the course to new contexts. There will be two midterm exams and a final exam.

\section*{Grading Scheme}
\subsection*{Component Weights}

\noindent\begin{tabular}{ll}
\textbf{Component} & \textbf{Weight}       \\ \hline
\textbf{Field work}             & 15\%  \\
\textbf{Homework}               & 15\%       \\
\textbf{Reflection}             &  8\%   \\
\textbf{Engagement}             &  7\%       \\
\textbf{Labs}                   & 20\% \\
\textbf{Exams}                  & 35\%
\end{tabular}%

\lettergrades{a) a pattern of seeking help in the course at office hours, b) number of missing homework assignments.}

\section*{Policies}
\subsection{AI usage}
\textbf{AI usage is strictly prohibited on reflection assignments and physics field work.} These assignments will be hand-written and require signing the following statement: ``I acknowledge that the use of generative AI on this assignment represents a cowardly avoidance of engagement with my own personal experience.'' 

\subsection{General course policies}
Expect a response within one business day for emails sent to the instructor. To ensure a prompt response, format the subject line as: [PHY 1200] *insert subject here*

Students are expected to arrive and leave on time. I will attempt to arrive five minutes early and be available for at least five minutes after class.
\latework{Partial}
\collegepolicies

\end{document}
